% Commentary: GPT-6 Astra (OpenAI, ChatGPT Pro mode).

\begin{dcomment}{context}{p1}{Why this construction belongs in a theory of derivators}
Grothendieck needs factorizations which can be used uniformly across a category of coefficients. A construction for isolated arrows would not supply the functorial replacements needed when those coefficients occur in diagrams. The emphasis on an endofunctor of the arrow category, and then on an entire transfinite sequence of such endofunctors, should be read with that purpose in view.

Two strands of the chapter have different origins. Its language of lifting and factorization is that of homotopical algebra; its treatment of size continues the cardinal filtrations and accessibility developed in \citep[Expos\'e~I, \S9]{SGA4I1972}. A filtration does more than assign a cardinal label to an object: it specifies presentations, closure under controlled colimits, and the bounds within which represented functors commute with those colimits. Here that apparatus will turn an unbounded repetition of a local operation into a uniform existence theorem.

The commentary on Chapter~VI already discusses the two forms of the replacement axiom, by a functorial replacement and by a factorization ``à la Quillen'', at manuscript p.~108, and the distinction between replacements at manuscript p.~180. The present pages address the additional existence problem: how to manufacture the replacements from a set of maps. They should therefore be read as preparation for the diagram theory, rather than as another formulation of its axioms. The modern account of the same small object mechanism in \citep[Chapter~12]{Riehl2014} makes this separation of construction and use especially clear.
\end{dcomment}

% Commentary and mathematical audit by GPT-6 Astra (OpenAI, ChatGPT Pro mode).
\begin{dcomment}{exposition}{p1}{Factorization by adjoining solutions}
Grothendieck wants a construction which can be used uniformly for every arrow and subsequently compared as its generators vary. The object to be constructed is therefore more than an isolated factorization. It is a section of the composition functor
\[
 \mathcal M^{[2]}\longrightarrow\mathcal M^{[1]},
 \qquad (X\longrightarrow E\longrightarrow Y)\longmapsto(X\longrightarrow Y),
\]
where $[n]=\{0<\cdots<n\}$. This explains the apparently leisurely passage through the arrow category: naturality in a commutative square will eventually supply natural replacements and comparisons between them.

Write $I=C_0$. In modern lifting notation the two desired classes are
\[
 R=I^\perp,\qquad L={}^{\perp}R.
\]
Here $j\perp p$ means that every commutative square with $j$ on the left and $p$ on the right admits a diagonal filler; uniqueness is not required. Thus Grothendieck's $(C_0)_*$ is $R$, and $\widetilde{C_0}$ is $L$. The construction will first produce a much more concrete subclass of $L$, consisting of maps built by attaching copies of the generators. Retracts enter only when one identifies all of $L$.

There are two independent tasks. Adjoining a cell for every lifting problem makes all problems present at one stage solvable at the next. Smallness of the domains then ensures that a problem appearing at a sufficiently long limit was already present earlier. Neither task asserts that one application solves every problem in its newly enlarged domain. This distinction is the reason for the transfinite sequence.

The resulting argument is the classical small object argument, whose modern formulation and retract conclusion are given in \citet[Theorem~12.2.2 and Corollary~12.2.4]{Riehl2014}. Grothendieck isolates a useful economy in its hypotheses: arbitrary coproducts in $\mathcal M$ are unnecessary. Pushouts of the generating arrows, together with filtered colimits, suffice to form all the amalgamated sums under the fixed source. The word ``accessible'' concerns objects as tested by their covariant representable functors; in the proof below, only factorization of maps from the generating domains through stages is used. The categorical accessibility assumptions later in this chapter are a way of obtaining a uniform bound for that property.
\end{dcomment}

\begin{dcomment}{exposition}{p5}{A single stage and the colimits it needs}
The first stage is arranged to solve every prescribed extension problem without choosing one problem in preference to another. For $f:X\to Y$, its indexing set can be written explicitly as
\[
 \Delta(I,f)=\coprod_{(j:A\to B)\in I}
 \{(a,b)\in\operatorname{Hom}(A,X)\times
                   \operatorname{Hom}(B,Y)\mid fa=bj\}.
\]
This is a small set: $I$ is small and the category is locally small in the chosen universe. Each triple $D=(j,a,b)$ gives the pushout $X_D=X\amalg_A B$, together with its induced map to $Y$. In the category of objects under $X$, take their coproduct. This gives exactly (1.23), and the inclusion of the $D$th cell supplies the map $B\to X_1(f)$ required by (1.17).

The existence argument deserves to be explicit because the stated assumptions do not include all small colimits. A pushout of a pushout of $j$ is again a pushout of $j$: this is the pasting property of pushout squares. Consequently every map $X\to X_D$ is itself pushoutable along every map out of $X$. Any finite family of these objects therefore has an amalgamated sum under $X$, built by repeated binary pushouts. For a finite subset $F\subseteq\Delta(I,f)$, write this sum as $X_F$, and set $X_{\varnothing}=X$. Chosen colimit cocones give transition maps for inclusions $F\subseteq G$. The poset of finite subsets is filtered, and
\[
 X_1(f)=\mathop{\mathrm{colim}}_{F\subseteq\Delta(I,f),\;F\text{ finite}} X_F.
\]
To verify its universal property, a map from this colimit to an object under $X$ is a compatible family of maps from the finite amalgamations. Restriction to singleton subsets identifies such a family with precisely one compatible map from each $X_D$. This proves the required coproduct property in the undercategory, including the case $\Delta(I,f)=\varnothing$.

If ordinary coproducts also exist, the same object is the pushout
\[
\begin{tikzcd}
 \coprod_{D\in\Delta(I,f)} A_D \arrow[r]\arrow[d] & X\arrow[d,"i_1(f)"]\\
 \coprod_{D\in\Delta(I,f)} B_D \arrow[r] & X_1(f).
\end{tikzcd}
\]
This familiar presentation is convenient for calculations; the preceding finite-subset construction explains why it is not an additional assumption in Theorem~1. The construction of a stage uses the codomains $B_D$, but the eventual lifting argument will require smallness only of the domains $A_D$.
\end{dcomment}

\begin{dcomment}{exposition}{p7}{Why the cofinality bound solves every problem}
Grothendieck's Corollary~2 supplies the decisive passage from a succession of partial solutions to an actual right lifting property. The statement needs the additional word \emph{limit}. For a limit ordinal, the condition that every subset of cardinal at most $\pi$ has an upper bound is equivalent to cofinality greater than $\pi$. More invariantly, choose a regular cardinal $\kappa$ for which every generating domain is $\kappa$-presentable, and use a limit ordinal $\lambda$ with $\operatorname{cf}(\lambda)\geq\kappa$. The ordered set of stages $\beta<\lambda$ is then $\kappa$-filtered.

Consider a square $j:A\to B$ against $p_\lambda:X_\lambda\to Y$, with top map $a:A\to X_\lambda$ and bottom map $b:B\to Y$. Since
\[
 X_\lambda=\mathop{\mathrm{colim}}_{\beta<\lambda}X_\beta,
\]
smallness gives $\beta<\lambda$ and $a_\beta:A\to X_\beta$ representing $a$. The compatibility of the maps to $Y$ gives
\[
 p_\beta a_\beta=p_\lambda a=bj.
\]
Thus the square already commutes at this stage. The parenthetical suggestion to enlarge $\beta$ is harmless but unnecessary: the target $Y$ has remained fixed throughout. The corresponding cell at stage $\beta+1$ provides $h:B\to X_{\beta+1}$ with $hj=t_{\beta+1,\beta}a_\beta$ and $p_{\beta+1}h=b$. Because $\lambda$ is a limit, $\beta+1<\lambda$, and $t_{\lambda,\beta+1}h$ is the desired filler.

No smallness assertion about $B$ was used. In fact the proof needs only the surjectivity of the map from the colimit of the relevant Hom sets, rather than its injectivity. This is why the usual relative-smallness hypothesis, restricted to the cellular chains actually constructed, is sufficient as well. The bound must be common to all generating domains, and the stopping ordinal must be common to all input arrows if one wants a single factorization functor. Since the generators form a set, individual accessibility bounds can be replaced by a common larger one.

The omitted limit condition is a real, although easily repaired, boundary error. In the definition used in \citet[Expos\'e~I, Definition~9.1]{SGA4I1972}, a poset with a maximum is ``grand devant $\pi$'' for every $\pi$. Thus $I_1=\{0\}$ satisfies the printed hypothesis. Take sets, generators $\varnothing\to 1$ and the fold $2\to 1$, and $f=1_1$. The first generator adds one point; the square for the fold identifies no distinct old points. Hence $X_1(f)=2$ and $p_1(f):2\to1$ is not injective. Since the right class for these two generators consists of bijections, Corollary~2 would be false at $\alpha=1$. Its proof itself shows the repair: choose a limit ordinal with the indicated cofinality, as in the successor-cardinal example given in the text. The same qualification is required in (1.22) and Theorem~1(b).

The two printed labels $p_{\alpha+1}$ at the end of the proof, on printed page~7, omit the prime on $\alpha$. Their intended value is $p_{\alpha'+1}$. The displayed source of that arrow determines the correction; the edition's readings are retained in the French text and registered in the corrections table.
\end{dcomment}

\begin{dcomment}{examples}{p8}{Empty families and a complete calculation in sets}
Grothendieck has now separated attachment from lifting, and the simplest generator makes both operations visible. Work in sets and let $I=\{\varnothing\to\{*\}\}$. For $f:X\to Y$, a lifting problem is exactly an element of $Y$. Hence
\[
 X_1(f)=X\amalg Y,\qquad i_1(f)=\mathrm{in}_X,
 \qquad p_1(f)=[f,1_Y].
\]
The right class is the class of surjections, because a lift against $\varnothing\to\{*\}$ is a preimage of one specified element. With the usual axiom of choice, the left class is the class of injections. An injection lifts against a surjection by extending the prescribed map over its complement. Conversely, if $g:A\to B$ is not injective, choose distinct $a,a'$ with $g(a)=g(a')$, map $A$ to $\{0,1\}$ with different values at these points, and use the surjection $\{0,1\}\to\{*\}$. A filler would factor that map through $g$, which is impossible.

The whole transfinite construction can be calculated, with the obvious compatible choices of coproducts:
\[
 X_\alpha(f)=X\amalg(Y\times\alpha),\qquad
 p_\alpha(x)=f(x),\qquad p_\alpha(y,\beta)=y.
\]
At a successor one adds another copy of $Y$; at a limit one takes the union of the earlier copies. For every $\alpha>0$, $p_\alpha$ is already a surjection, but if $Y\ne\varnothing$ every transition to the next stage adds elements. Success in the lifting problem therefore does not mean that the unmodified construction has stabilized.

There is also a small literal omission in (1.18)--(1.20). The definition of $C_0''$ expressly allows only \emph{nonempty} amalgamated families. In this example $C_0'$ consists of the injections which add one element, and $C_0''$ consists of the injections with nonempty complement. For $f:\varnothing\to\varnothing$, however, $\Delta(I,f)$ is empty and (1.23) gives $i_1(f)=1_{\varnothing}$. Thus (1.20), read with the printed definition, is false in this boundary case. The minimal repair is to permit the empty amalgamated family, whose structural map is the identity, or explicitly to adjoin identities to $C_0''$. Under the printed convention the always valid assertion is $i_1(f)\in C_0''\cup\mathrm{Iso}(\mathcal M)$. This is a convention repair in the source, not a transcription correction, and it does not affect the resulting weak factorization theorem. Zero-length cellular compositions must in any event include identities.
\end{dcomment}

\begin{dcomment}{exposition}{p9:Il en résulte aussitôt, par récurrence ordinale}{The cellular closure and its successive steps}
The finer class $C_0'''$ records how a map was constructed, before retracts erase that information. Its intended modern description is the class of relative $I$-cell maps: continuous transfinite composites in which every successor map is a pushout of an amalgamated family of copies of generators. Identities are allowed, and all diagrams are small in the chosen universe. The separate stages $C_0'$ and $C_0''$ distinguish one attachment from simultaneous attachment of a family.

The assertion that these maps lie in ${}^{\perp}(I^\perp)$ has a direct proof. Fix $q:U\to V$ in $I^\perp$. A lifting problem for a pushout of $j\in I$ reduces, by the pushout property, to one for $j$. For an amalgamated family of such maps, choose a filler for each member with the common prescribed value on the source; the resulting maps agree there and therefore induce a filler on the amalgamation. For a continuous chain $(Z_\beta)_{\beta\leq\lambda}$, construct the lift on $Z_\beta$ by transfinite recursion. At a successor extend the lift already constructed using the lifting property of $Z_\beta\to Z_{\beta+1}$. At a limit the compatible lifts determine a map from the colimit. This proves the lifting property for $Z_0\to Z_\lambda$. It uses existence of fillers, together with the ordinary choice principle, and no uniqueness.

The word ``successor'' is essential in this formulation. In the explanatory sentence of (1.19), the displayed symbols $i_{\beta,\alpha}$ can be read as demanding that \emph{every} transition already belong to the class. That condition by itself would not imply closure under finite composition. For example, in the poset category $0<1<2$, take $I=\{0\to1,1\to2\}$. Pushouts and amalgamations give these two arrows and identities, but not $0\to2$. A chain all of whose transitions lie in this class cannot contain both $0$ and $2$; its colimit introduces no missing arrow. Nevertheless two successive cell attachments factor $0\to2$ through $1$. Thus the induction in (1.21) requires the usual successorwise interpretation of ``limites inductives ordinales strictes'', or an explicit additional closure under composition. We adopt that intended interpretation throughout; the claim about the lifting class is valid independently of this notational compression.
\end{dcomment}

\begin{dcomment}{exposition}{p10:Il faudrait quand même expliciter la dépendance fonctorielle}{The omitted functoriality in the input arrow}
The footnote asks for the functorial dependence which makes Theorem~1 usable. Its proof is short once the indexing squares are retained as data. A morphism $(u,v):f\to f'$ in the arrow category consists of $u:X\to X'$ and $v:Y\to Y'$ with $vf=f'u$. It sends
\[
 (j,a,b)\longmapsto(j,ua,vb)
\]
to a well-defined member of $\Delta(I,f')$. For this image $D'$, the pair of maps $X\xrightarrow{u}X'$ and $1_B$ induces a map
\[
 X_D=X\amalg_A B\longrightarrow X'\amalg_A B=X'_{D'}.
\]
These maps respect the structural maps from $X$ and the maps to the target. Composing them with the inclusions into $X_1(f')$ gives a cocone over the family defining $X_1(f)$, hence a unique map $X_1(u,v):X_1(f)\to X_1(f')$.

Identity and composition laws follow by restricting to $X$ and to each copy of $B$: there the asserted equalities are the identity and composition laws for $u$ and $v$, and the universal property makes their extensions unique. The same argument proves the naturality of $i_1$ and $p_1$. At successor ordinals one applies this functor to the already constructed square between $p_\alpha(f)$ and $p_\alpha(f')$; at limits one takes the map of colimits. Transfinite induction gives all of (1.10)--(1.14), including compatibility with every transition map.

One must fix the relevant colimit cocones to obtain literal functors. Another choice gives canonically isomorphic functors by the same universal properties. There is no need to choose fillers in arbitrary lifting problems: the attached cell itself supplies the particular solution required at each stage. This distinction explains why the factorization is functorial even when the fillers subsequently used to characterize the left or right class are highly nonunique.
\end{dcomment}

\begin{dcomment}{exposition}{p12:Donc si $\Phi_0$ est stable par facteurs directs}{Why the retract conclusion is exact}
The corollary identifies the amount of closure required after the cellular construction. Put $R=I^\perp$ and $L={}^{\perp}R$. If $I\subseteq\Phi_0\subseteq L$, then $\Phi_0^\perp=R$: the inclusion $I\subseteq\Phi_0$ gives one containment, and $\Phi_0\subseteq{}^{\perp}R$ gives the other. Assume, as in the corollary, that every map factors as a map in $\Phi_0$ followed by a map in $R$.

Take $g:A\to B$ in $L$ and choose a factorization
\[
 A\xrightarrow{j}E\xrightarrow{q}B,
 \qquad j\in\Phi_0,\quad q\in R,\quad qj=g.
\]
The square with left map $g$, top map $j$, right map $q$ and bottom map $1_B$ admits a filler $s:B\to E$. It satisfies $sg=j$ and $qs=1_B$. In the arrow category the maps
\[
 g\xrightarrow{(1_A,s)}j\xrightarrow{(1_A,q)}g
\]
compose to the identity. Thus $g$ is a retract of a member of $\Phi_0$, with the source held fixed. Conversely a retract of a map with the left lifting property has that property: insert the retract maps into a lifting square, solve the resulting square, and compose back. Therefore $L$ is exactly the retract closure of $\Phi_0$.

The cellular class and its retract closure consequently have different jobs. Cellular maps are the output of the explicit construction; the retract closure is the class characterized purely by lifting. In particular, if $\Phi_0$ is retract closed then it is already $L$. The stronger description using $C_0'''$ follows from the transfinite attachment construction, with identities included as explained above. This is the content, rather than an additional hypothesis, of the assertion that $(\Phi_0,R)$ is a Quillen precouple. Section~4 will isolate the same retract mechanism without any reference to the particular generating set.
\end{dcomment}

\begin{dcomment}{later}{p12:Donc si $\Phi_0$ est stable par facteurs directs}{From a pointed endofunctor to algebraic factorization}
Grothendieck's emphasis on $1\to p_1$ makes the construction especially close to the later algebraic treatment. It should not, however, be read as already providing a monad: a pointed endofunctor comes with a unit but has no specified multiplication or associativity law. The set example above exhibits the difference concretely, since the ordinary iteration keeps adding solutions to problems which have already been solved.

\citet[Theorem~4.4]{Garner2009} constructs an algebraically free natural weak factorization system on a small category of generating arrows under specified cocompleteness and smallness hypotheses, in particular for a locally presentable category. Its right factor carries a monad and its left factor a comonad. For a discrete generating category, the algebras on a right map encode chosen solutions of the generating lifting problems; morphisms preserve those choices. The added algebraic structure gives a universal property to the construction. It is stronger than the existence of the functorial factorization obtained here.

This places the chapter's two concerns in their subsequent setting. The transfinite argument proves that the required lifting classes form a weak factorization system, while the algebraic refinement controls the choices of solutions and their compatibility. A functorial section of composition, an algebraic weak factorization system, and an orthogonal factorization system with unique fillers are three distinct structures. Nothing in Theorem~1 requires unique fillers.
\end{dcomment}

\begin{dcomment}{exposition}{p15:commutativité qui revient à dire que}{Functoriality in the generators, including limit stages}
Grothendieck next keeps track of what happens when more lifting problems are declared relevant. The one-stage comparison is induced by the inclusion of indexing sets, but at later stages the input to the attachment functor has also changed. Writing the recursive formula makes the joint naturality precise.

Let $I\subseteq J$, put $P_I=p_1^I$ and $P_J=p_1^J$, and let $\tau:P_I\to P_J$ be the natural transformation constructed from the inclusion of squares. Both endofunctors lie over the codomain functor, and their units satisfy $\tau\eta_I=\eta_J$. Let $u_\alpha:P_I^\alpha\to P_J^\alpha$ denote the comparison at stage $\alpha$, where the superscript means the continuous iteration used in the text, not ordinary iteration at a limit. Then
\[
 u_0=1,\qquad
 u_{\alpha+1}=P_J(u_\alpha)\,\tau_{P_I^\alpha},\qquad
 u_\lambda=\mathop{\mathrm{colim}}_{\alpha<\lambda}u_\alpha
 \quad(\lambda\text{ a limit}).
\]
The successor formula starts with $P_I(P_I^\alpha f)$, changes the outer attachment functor to $P_J$, and then changes its input to $P_J^\alpha f$. Naturality of $\tau$ allows these two operations in the opposite order as well.

Here is the compatibility with a successor transition, which is the point needed to make the limit formula meaningful. Naturality of $\eta_J$ and the equation $\tau\eta_I=\eta_J$ give
\[
 \begin{split}
 u_{\alpha+1}\eta_{I,P_I^\alpha}
 &=P_J(u_\alpha)\eta_{J,P_I^\alpha}\\
 &=\eta_{J,P_J^\alpha}u_\alpha.
 \end{split}
\]
Compatibility with all transitions follows by composition and the defining colimit property at limits. Each $u_\alpha$ is natural in $f$ by the same induction. For $I\subseteq J\subseteq K$, the equation $u_\alpha^{K,I}=u_\alpha^{K,J}u_\alpha^{J,I}$ holds at stage one because both maps restrict to the same inclusions of cells. At successors it follows by naturality of the one-stage transformations, and at limits by equality on every colimit component. Identity inclusions are handled identically. This proves the functoriality on the product of the poset of generating sets with the arrow category that the final paragraph of Section~2 leaves unexpanded.

The term ``inclusion'' describes the indexing sets. It does not by itself assert that the resulting morphism of objects is a monomorphism in an arbitrary category. The following elementary refinement explains when the same word can also be used for the actual stages.
\end{dcomment}

\begin{dcomment}{development}{p16:je n'explicite pas le détail de la fonctorialité}{Comparisons which are monomorphisms}
The parameter dependence in Section~2 is useful when a family of generators is enlarged gradually. For presheaves, monomorphic generators give a stronger conclusion: an enlargement can be seen stage by stage inside the resulting objects. We record the precise two-variable statement, with its proof, for this construction. We checked the small object construction in \citet[Section~12.2]{Riehl2014} and the algebraic comparison in \citet[Sections~3--4]{Garner2009}; we did not find this particular formulation there. The statement below is an elementary elaboration of the displayed construction, with no claim of priority.

\textbf{Proposition (monomorphic comparison).}\ours
Let $\mathcal E=\mathbf{Set}^{\mathcal A^{\mathrm{op}}}$ for a small category $\mathcal A$. Let $I\subseteq J$ be small sets of monomorphisms in $\mathcal E$. Suppose $(u,v):f\to f'$ is a commutative square whose source map $u:X\to X'$ and target map $v:Y\to Y'$ are both monomorphisms. For every ordinal $\alpha$ for which the construction is performed, the jointly natural comparison
\[
 X_\alpha^I(f)\longrightarrow X_\alpha^J(f')
\]
is a monomorphism. In particular, at a fixed input $f$, enlargement of a set of monomorphic generators gives monomorphisms at every stage.

\textit{Proof.} First consider one stage. The map
\[
 \Delta(I,f)\longrightarrow\Delta(J,f'),\qquad
 (j,a,b)\longmapsto(j,ua,vb)
\]
is injective. Indeed equal images have the same generator $j$, and cancellation of the monomorphisms $u$ and $v$ gives equality of both maps $a$ and $b$. Evaluate the pushout formula at an object $c\in\mathcal A$. Evaluation preserves colimits, and every $j:A_j\to B_j$ is pointwise injective. Thus, as sets,
\[
 X_1^I(f)(c)\cong X(c)\amalg
 \coprod_{(j,a,b)\in\Delta(I,f)}
 \bigl(B_j(c)\setminus j(A_j(c))\bigr).
\]
This decomposition is only used at the level of sets; the complements need not form subpresheaves. On the first summand the comparison is $u(c)$, which is injective. On each remaining summand it is the identity on the same complement, placed in the summand with image index. Injectivity of the indexing map shows that distinct summands remain distinct. Hence the one-stage comparison is pointwise injective.

At stage zero the assertion is the hypothesis on $u$. If it holds at stage $\alpha$, apply the one-stage result to the square between $p_\alpha^I(f)$ and $p_\alpha^J(f')$; its source map is monic by induction and its target map is still $v$. This proves the successor step. At a limit $\lambda$, evaluate at $c$ and use the two compatible diagrams of sets. If two representatives in the first colimit have equal images in the second, they become equal at a common later stage of the second diagram. The injectivity of the comparison at that stage makes their images in the first diagram equal there as well. They therefore represent the same element of the first colimit. This proves the limit step and the proposition.\hfill$\square$

The assumption on the target map carries real information. With $I=J=\{\varnothing\to\{*\}\}$, take $f:\varnothing\to\{0,1\}$ and $f':\varnothing\to\{*\}$, the identity on the empty source and the constant map on targets. The induced map between first stages is $\{0,1\}\to\{*\}$, so it is not monic. The construction distinguishes lifting problems by their target component as well as their boundary map.
\end{dcomment}


% Commentary and audit: GPT-6 Astra (OpenAI, ChatGPT Pro mode).
\begin{dcomment}{exposition}{p16}{A filtration measures the cost of a construction}
Grothendieck now asks for more than the existence of a factorization. He wants to perform the construction inside a bounded part of the category, uniformly for every arrow in that part. This is what will eventually let a small amount of data control a model structure. The three quantities to keep separate are the size of the domains and codomains of the generators, the number of generators actually used, and the length of the construction. Isomorphic generators have the same lifting property, but the unmodified construction attaches a cell for each of them.

The terminology comes from \citet[Expos\'e I, Definitions 9.1--9.3 and 9.12]{SGA4I1972}. A partially ordered set is \emph{large relative to} an infinite cardinal $\pi$ if each of its subsets of cardinality at most $\pi$ has an upper bound. Thus the modern directedness index is $\pi^+$. An object called $\pi$-accessible here has a representable functor preserving such colimits; in the usual setting this is $\pi^+$-presentability, not presentability at the cardinal $\pi$. The filtration is additional structure: its levels are essentially small, are increasing, admit bounded presentations of objects at larger levels, and have specified closure properties. Membership in a filtration level should therefore not be replaced by an informal assertion that an object has at most $\pi$ elements.

For the calculation ahead we need the following concrete consequences. Put $\mathcal M^\pi=\Filt^\pi$. If $A\in\mathcal M^r$, it has a filtered presentation by objects of $\mathcal M^{\pi_0}$ with at most $r^{\pi_0}$ indices. If $X\in\mathcal M^s$ and $s\geq\pi_1$, it has a presentation by objects of $\mathcal M^{\pi_1}$ over a $\pi_1^+$-directed poset with at most $s^{\pi_1}$ indices. These are the presentations used on ms.~20. The stronger closure condition (3.13), which is explicitly assumed below, includes \emph{all} filtered diagrams of cardinality at most the level, not only those large relative to $\pi_0$.

This stronger closure also supplies the retract argument in (3.2). If $X$ is a retract of $Y\in\mathcal M^\pi$, with idempotent $e:Y\to Y$, the countable chain
\[
Y\xrightarrow{e}Y\xrightarrow{e}Y\xrightarrow{e}\cdots
\]
has colimit $X$, so $X\in\mathcal M^\pi$. If $X$ is $\pi$-accessible, factor its identity through a stage of its $\pi^+$-directed presentation by $\mathcal M^\pi$-objects. This makes it such a retract. In the other direction, a presentation with at most $\pi^{\pi_0}$ arrows expresses $\Hom(X,-)$ as a limit of that many representables. Directed colimits of sets with index larger than this bound commute with those limits. This proves the inclusions in (3.2) under the closure hypotheses used here and explains why their equality at suitable cardinals is useful.
\end{dcomment}

\begin{dcomment}{exposition}{p17}{The cofinality qualification in the cardinal arithmetic}
Grothendieck looks for cardinals at which the estimate stops growing. The unconditional choice $\pi=2^c$, $c\geq\pi_0$, works because
\[
(2^c)^{\pi_0}=2^{c\pi_0}=2^c.
\]
The parenthetical appeal to the generalized continuum hypothesis has a different status. Even under GCH, $\pi>\pi_0$ does not imply $\pi^{\pi_0}=\pi$. For example, $\pi=\aleph_\omega$ and $\pi_0=\aleph_0$ give $\pi^{\pi_0}>\pi$. More generally, K\"onig's theorem gives $\pi^{\operatorname{cf}\pi}>\pi$ for every infinite $\pi$. A direct proof of the case just used is instructive: if $(f_\xi)_{\xi<\aleph_\omega}$ enumerated all maps $\mathbb N\to\aleph_\omega$, choose $g(n)$ outside the set $\{f_\xi(n):\xi<\aleph_n\}$. For every $\xi$ some $n$ satisfies $\xi<\aleph_n$, and then $g(n)\ne f_\xi(n)$, a contradiction.

Under GCH the additional condition $\pi_0<\operatorname{cf}\pi$ is sufficient. The general exponentiation rule and its cofinality qualification are recorded in \citet[Theorems 10.12 and 11.7, Exercise 11.8]{Marks2026}. Nothing in the factorization theorem requires GCH: the fixed exponent identity above is all that is needed. The same correction applies to the repeated GCH assertion on ms.~19.

The sentence about the arrow category immediately before (3.4) should also be read with its size threshold in view. Pointwise filtered colimits and essential smallness of the bounded arrow categories are immediate. Uniformly bounded presentations of arrows require control of the Hom sets between small objects and simultaneous realization of the arrow in two presentations. At a sufficiently enlarged initial level this is the usual accessibility theorem for a finite diagram category. The text does not supply that extra argument at the original $\pi_0$, and does not use the assertion in its proof.
\end{dcomment}

\begin{dcomment}{examples}{p19}{An orbit count is not a bound on a Hom set}
The purpose of $C_1$ is to pay for every choice of a map between the small objects. As printed, however, it counts isomorphism classes of objects of the arrow category. These two numbers can be very different. This is a mathematical issue in the edition, not an error introduced by the transcription.

Let $G$ be any small group and let $BG$ have one object $*$, with $\Hom(*,*)=G$. Give $BG$ the constant filtration $\mathcal M^\pi=BG$, with $\pi_0=\pi_1=\aleph_0$. Every filtered diagram in $BG$ is isomorphic to the constant diagram: transport from a fixed index to another through a common upper bound; a further common upper bound shows that transport does not depend on the choices. Consequently filtered colimits exist, and $\Hom(*,-)$ preserves them. Pushouts of spans of isomorphisms exist as well. The constant filtration has all the presentation and closure properties needed here.

Every two arrows $g,h$ are isomorphic as objects of $\Fl(BG)$: take the square with vertical maps $1$ and $hg^{-1}$. Thus the printed constant is $C_1=1$, whereas $|\Hom(*,*)|=|G|$. Take $r=s=\aleph_0$ and a group of cardinality greater than $\aleph_0^{(\aleph_0^{\aleph_0})}$. Formula (3.8), with the printed $C_1$, is false. The example also identifies precisely the invalid step on ms.~20: the indicated bound $|\Hom(A_i,X_j)|\leq C_1$ does not follow from the definition of $C_1$.

There is a straightforward repair. Choose a small full skeleton $K_1$ of $\mathcal M^{\pi_1}$, choosing its representatives compatibly with those at level $\pi_0$, and use
\[
\chi=\max\{\aleph_0,\pi_1,|\mathrm{Mor}(K_1)|\}
\]
as the threshold for $r$ and $s$. It counts actual maps between representatives, including automorphisms. This change is also the natural reading of the lowercase $c_1$ in Theorem~2: the entire strictly full level need not be a small category, so its raw collection of arrows cannot serve as a small cardinal. The choice of skeleton affects a numerical bound, not the factorization or the lifting classes.
\end{dcomment}

\begin{dcomment}{exposition}{p22}{The estimates survive the corrected counting argument}
Grothendieck's desired bound has enough room to absorb the choices that were omitted in the displayed counting. We give the calculation with the corrected threshold $r,s\geq\chi$ of the preceding block. Write $\mu=r^{\pi_0}$. For the presentations on ms.~20, accessibility gives
\[
\Hom(A_i,X)\cong\varinjlim_j\Hom(A_i,X_j),
\qquad |\Hom(A_i,X)|\leq s^{\pi_1}.
\]
Taking a limit over the presentation of $A$ then yields
\[
|\Hom(A,X)|\leq (s^{\pi_1})^{|I|}
\leq s^{\pi_1\mu}=s^\mu.
\]
This proves the repaired (3.8). The last equality uses $\pi_1\leq r\leq\mu$ and infinite cardinal arithmetic.

The choices of all the objects $A_i$ must be counted by a product, not by $|I|$ times the number of possibilities at one index. The same applies to the transition arrows. A poset with at most $\mu$ objects has at most $\mu$ arrows; a diagram on it with values in a fixed small skeleton at level $\pi_0$ can therefore be specified in at most $\chi^\mu\leq 2^\mu$ ways. The requirements of identity and composition merely reduce this number. There are at most $2^\mu$ isomorphism types of posets of size at most $\mu$, since both their underlying sets and their order relations can be coded inside a fixed set of size $\mu$ and its square. Taking colimits gives
\[
|\operatorname{IsoOb}(\mathcal M^r)|\leq 2^\mu.
\]
This replaces the unsupported fixed-index bound $r^{\pi_0}$ on ms.~21 and~23 by $2^\mu$; the eventual bound is unchanged.

For a fixed $X\in\mathcal M^s$, choose one representative of each object at level $r$. There are at most $s^\mu$ arrows from any such representative to $X$, hence at most $2^\mu s^\mu=s^\mu$ objects of $\mathcal M^r/X$ up to isomorphism. Between two of them there are at most $r^\mu=2^\mu$ morphisms, because a morphism in the slice is a particular underlying morphism in $\mathcal M$. Consequently
\[
|\operatorname{IsoOb}(\Fl(\mathcal M^r/X))|
\leq(s^\mu)^2 2^\mu=s^\mu.
\]
This is (3.9), including the case $s<r$. Finally, applying the same argument to the arrow category without a fixed target gives (3.7). In particular the appearance of products where exponential counts are required is a false intermediate inference, but it does not force a worse final estimate once the small skeleton has been counted correctly.

The elementary identity used here is $a^b=2^b$ for infinite $b$ and $2\leq a\leq b$: one has $2^b\leq a^b\leq(2^b)^b=2^b$. These are the hypotheses under which Lemma~3.2 is used. Its unqualified wording should not be extended to arbitrary finite cardinals, for which, for example, $3^3\ne2^3$.
\end{dcomment}

\begin{dcomment}{examples}{p23}{What the proposed sharpness example does show}
Grothendieck rightly observes that the Hom bound cannot improve on the elementary set-theoretic exponent $s^r$: sets of cardinalities $r$ and $s$ have exactly that many functions. The claim that the category of pairs $(A,B)$, $B\subseteq A$, attains the arrow-isomorphism bound $2^r$ at every allowed $r$ needs a separate check. First, the printed inequality $|A|\geq\pi$ on ms.~20 must be $|A|\leq\pi$ to describe the intended increasing bounded filtration; the original sign remains in the source text and is recorded in the corrections register.

Even after that change, equality does not hold at every cardinal. Assume GCH and take $r=\aleph_1$, $\pi_0=\pi_1=\aleph_0$. A map of such pairs is a function between sets with two colours, required to preserve the distinguished colour. Its isomorphism type is determined by the colour of a target element, the two cardinalities of its source fibre, and the multiplicity of each such fibre type. Among cardinalities at most $\aleph_1$ there are only countably many choices. Thus there are countably many fibre types, and at most
\[
\aleph_1^{\aleph_0}=\aleph_1
\]
possible multiplicity lists. The count of arrow-isomorphism types is therefore at most $\aleph_1$, strictly less than $2^{\aleph_1}=\aleph_2$. The threshold involving the printed $C_1$ is met: arrows between countable coloured sets have at most $2^{\aleph_0}=\aleph_1$ isomorphism types.

The same example does attain $2^{\aleph_0}$ at the countable level. For each subset $S$ of the positive integers, take one fibre of size $n$ for each $n\in S$, and add a countably infinite fibre so that the source remains countable. Distinct $S$ give nonisomorphic arrows. This is an orbit-count example at a level below the proposition's threshold $C_1=2^{\aleph_0}$, rather than an admissible equality case for that proposition. These calculations retain the useful point of the example while separating a genuine sharp Hom estimate from a stronger, false, universal sharpness assertion about arrow orbits.
\end{dcomment}

\begin{dcomment}{exposition}{p26}{One fixed exponent controls every stage}
The geometric calculation has now become very simple. If $C_0\subseteq\Fl(\mathcal M^r)$ and there are at most $2^\mu$ generators, where $\mu=r^{\pi_0}$ and $r\geq\chi$, then for an arrow with endpoints in $\mathcal M^s$, $s\geq r$, the set of attaching squares has cardinality at most $s^\mu$. Each individual pushout stays in $\mathcal M^s$. Their common-source amalgam can be obtained as a filtered colimit of finite partial amalgams, using at most $s^\mu$ indices. This proves Lemma~3.3 with the corrected counting constant. If there are no squares, the amalgam is simply the original source; this case is included by allowing the empty finite subset, or by treating it separately.

What matters for iteration is an exponent fixed point. If $S=c^t$ with $c\geq2$ and $t\geq\mu$, then $S^\mu=c^{t\mu}=c^t=S$. Thus the successor construction preserves $\mathcal M^S$. A limit stage indexed by an ordinal of cardinality at most $S$ stays there by (3.13). Transfinite induction proves that every stage up to any ordinal $\alpha$ with $|\alpha|\leq S$ has its middle object in $\mathcal M^S$.

To obtain the lifting property at the final stage, the ordinal must also be a limit of cofinality greater than $\mu$. The initial ordinal of $\mu^+$ is a uniform choice: $\mu^+\leq2^\mu\leq S$. The assertion in Theorem~2 that the initial ordinal of \emph{any} cardinal $r'>\mu$ will do is too broad as a justification of the cofinality condition; a singular $r'$ may have cofinality at most $\mu$. Choosing the successor cardinal, as the theorem itself suggests, repairs this point. The size estimate and the cofinality test perform different jobs and must both be retained.
\end{dcomment}

\begin{dcomment}{development}{p27}{Keeping the multiplicity of generators in the estimate}
The non-plethoricity hypothesis is convenient when one is free to choose a generating set. In a family of constructions, however, discarding isomorphic copies can obscure the given maps between generating data. We can keep those copies and make their cost explicit. The following is a quantitative completion of the calculation, with no claim of priority. We checked the cardinal-filtration argument in SGA~4 and the accessibility bounds for functorial factorizations in \citet[Theorem 4.4]{Low2016Heart}; neither is cited as asserting this exact formula with the raw multiplicity parameter.

\textbf{A uniform bound with multiplicities.}\ours
Assume the filtered-colimit, pushout and cardinal-filtration hypotheses of Section~3, including (3.13). Choose $\chi$ by counting the actual arrows of a small full skeleton of $\mathcal M^{\pi_1}$. Let $C_0$ be any small set of coquarrable arrows whose endpoints lie in $\mathcal M^r$, where $r\geq\chi$. Let $f$ have endpoints in $\mathcal M^s$, and let $\alpha$ be any small ordinal. Put
\[
\mu=r^{\pi_0},\qquad
q=\max\{2,r,s,|C_0|,|\alpha|\},\qquad B=q^\mu.
\]
Then $X_\beta(C_0,f)\in\mathcal M^B$ for every $\beta\leq\alpha$. If $\alpha$ is a limit with $\operatorname{cf}(\alpha)>\mu$, the terminal factor $p_\alpha(f)$ has the right lifting property against $C_0$.

\emph{Proof.} The cardinal $\mu$ is infinite, $B\geq q\geq r,s$, and
\[
B^\mu=q^{\mu\mu}=q^\mu=B.
\]
At stage zero the assertion follows from monotonicity of the filtration. Suppose it is true at $\beta$. The repaired Hom estimate applies to each endpoint of a generator and to each endpoint of $p_\beta(f)$, and gives at most $B^\mu=B$ possibilities for either side of an attaching square. There are therefore at most $|C_0|B^2\leq B$ squares. Each of their pushouts belongs to $\mathcal M^B$. A finite partial amalgam also belongs to $\mathcal M^B$, by induction on the number of its terms; the empty one is $X_\beta$. The poset of finite subsets of the set of squares is filtered and has cardinality at most $B$. Condition (3.13) puts its colimit, which is $X_{\beta+1}$, in the same level.

If $0<\lambda\leq\alpha$ is a limit, its diagram has cardinality $|\lambda|\leq|\alpha|\leq B$. Again (3.13) puts $X_\lambda=\varinjlim_{\beta<\lambda}X_\beta$ in $\mathcal M^B$. This completes the induction. Finally, a generator domain is $\mu$-accessible by (3.2), so any map from it to $X_\alpha$ factors through some $X_\beta$, $\beta<\alpha$. The prescribed map from the generator codomain to the fixed target then determines a square at that stage, solved at $\beta+1<\alpha$. Its image in $X_\alpha$ is the required lift.\hfill$\square$

In particular take $\alpha$ to be the initial ordinal of $\mu^+$. The formula gives an explicit answer to the size question (3.6) for the unaltered generating set. The extra parameter $|C_0|$ is necessary in this formulation: in $\mathrm{Set}$, arbitrarily many distinct copies of the isomorphism class $\varnothing\to\{*\}$ attach arbitrarily many points to $\varnothing\to\{*\}$ at the first stage, while the filtration levels of the generator endpoints remain fixed.
\end{dcomment}

\begin{dcomment}{later}{p28:Tout a été prouvé, si on se rappelle}{From bounded factorizations to small model structures}
Theorem~2 addresses a question that remains useful long after the existence of the small object argument is understood: can a chosen factorization be performed inside a controlled small part of the ambient category? In the locally presentable setting, \citet[Theorem 4.4]{Low2016Heart} gives a functorial weak factorization system whose functors preserve specified presentable objects under explicit bounds on the generators, their Hom sets and two regular cardinals. His Theorem~5.9 then tests when two such systems form a model structure by restricting them to an appropriate small full subcategory; Proposition~5.12 gives the corresponding statement for an arbitrary combinatorial model category.

These results use regular presentability indices and the relation of one cardinal being sharply below another. The numerical dictionary is therefore not obtained by replacing every $\Filt^\pi$ with the modern class of $\pi$-presentable objects. What is common is the argument's architecture: control one stage, choose a cardinal stable under its operations, and keep the entire factorization within that bound. This is the immediate background for the model-structure criteria that Grothendieck develops in Section~4.
\end{dcomment}


% Commentary and audit by GPT-6 Astra (OpenAI, ChatGPT Pro mode).

\begin{dcomment}{exposition}{p28}{From a factorization procedure to its two lifting classes}
Grothendieck has constructed factorizations from a set of generators. He now asks which parts of that construction survive when the generators and their cardinal estimates are forgotten. This change of viewpoint explains the apparently elementary opening of Section~4: the next existence argument will need to recognize a model structure from two factorization procedures, and a factorization procedure does not by itself determine the classes of maps which ought to be called cofibrations and fibrations.

Write $L^{\square}$ for the maps having the right lifting property against every member of $L$, and ${}^{\square}R$ for the corresponding left lifting class. Grothendieck's notation is $L_*=L^{\square}$ and $R^*={}^{\square}R$. A \emph{pr\'ecouple} supplies factorizations and the inclusion $L\subset{}^{\square}R$. A \emph{couple} supplies the equalities
\[
 L={}^{\square}R,\qquad R=L^{\square}.
\]
Thus a couple is what is now called a weak factorization system. Neither uniqueness of the diagonal filler nor a functorial choice of factorization belongs to this definition. The constructions of Section~1 provide functoriality under their additional smallness hypotheses. Modern definitions and the retract argument can be compared directly in \citet[Definition~11.2.1 and Lemma~11.2.3]{Riehl2014}.

The word \emph{clos} concerns equality with a lifting class. It must be distinguished from closure under the operations listed later in (4.4.3). Those operations preserve a left lifting class, but the converse requires an existence theorem. Section~1 furnishes that converse for classes generated by a small set of sufficiently small maps. This distinction will become decisive when the class $C\cap W$ is known to have the expected stability properties but no generating set has yet been found.

The expression ``facteur direct'' also has a categorical meaning here. It is a retract in the arrow category, without any additive direct-sum decomposition. For arrows $j:A\to B$ and $i:A'\to B'$, it means that there are morphisms of arrows $j\to i\to j$ whose composite is the identity of $j$. Proposition~4.2 shows that a factorization of $j$ lets us keep its source fixed. That stronger form is why one diagonal filler is enough to turn a lifting condition into a retract statement.
\end{dcomment}

\begin{dcomment}{exposition}{p30}{The retraction which completes both classes}
The role of Proposition~4.2 is to identify exactly what can be missing from a pr\'ecouple. Suppose $j:A\to B$ lifts against $R$, and choose a factorization $j=pi$ with $i:A\to Z$ in $L$ and $p:Z\to B$ in $R$. Apply the lifting property to
\[
\begin{tikzcd}[row sep=large,column sep=large]
A \arrow[r,"i"] \arrow[d,"j"'] & Z \arrow[d,"p"]\\
B \arrow[r,"\id_B"'] \arrow[ur,dashed,"q"] & B .
\end{tikzcd}
\]
The equations are $qj=i$ and $pq=\id_B$. Hence $(\id_A,q):j\to i$ and $(\id_A,p):i\to j$ exhibit $j$ as a retract of $i$. Conversely, a lifting property passes to retracts: extend a proposed lifting square along the section, solve it for the larger arrow, and compose the resulting filler with the retraction. These two observations prove
\[
{}^{\square}R=\operatorname{Retr}(L),\qquad
L^{\square}=\operatorname{Retr}(R)
\]
whenever $(L,R)$ is a pr\'ecouple. The dual assertion uses a factorization with target fixed.

This proof requires a factorization only of the particular arrow $j$ under consideration, as footnote~9 correctly emphasizes. There is no induction hidden in the argument and no completeness hypothesis on the category. Once all arrows have factorizations, closing the two classes under retracts therefore produces a weak factorization system without changing any chosen factorization. This is the precise content of the two corollaries on manuscript page~31.

It also explains a useful economy in model-category arguments. One first constructs maps in classes with convenient explicit descriptions, such as relative cell complexes. One then enlarges those classes by retracts. The enlargement does not require repeating the construction: lifting supplies the necessary retraction. The enlargement is nevertheless real, and omitting it from a recognition statement can change that statement's truth value.
\end{dcomment}

\begin{dcomment}{exposition}{p34}{Recognition by two compatible factorizations}
Grothendieck wants to describe weak equivalences through their two possible factorizations. The names $TC$ and $TF$ are provisional at this stage: the assumptions only give $TC\subset C\cap W$ and $TF\subset F\cap W$. Proposition~4.3 measures what must be added before those inclusions become the desired equalities.

The proof has two different mechanisms. A factorization
\[
j=pi,\qquad j\in C\cap W,\quad
i\in\widetilde{TC},\quad p\in\overline{TF}
\]
exhibits $j$ as a retract of $i$, because $C$ lifts against $\overline{TF}=C_*$. This gives $C\cap W\subset\widetilde{TC}$ without using a property of $W$ beyond membership of $j$ in it. In the other direction, factor an arbitrary $w\in W$ using $(C,TF)$. If $w=pi$, then $p\in TF\subset W$; cancellation in $W$ forces $i\in C\cap W$. This is the sole use of the two-out-of-three axiom in that half of the argument. The dual proof exchanges the two factorizations.

Consequently, the equivalences among (a), (a$'$), (b), (b$'$), and (c) need only the restricted cancellation condition stated in footnote~10: whenever $gf\in W$, membership of $f$ in $W$ is equivalent to membership of $g$ in $W$. A displayed equality $W=TF\circ TC$ additionally uses closure of $W$ under composition for the reverse inclusion. The source itself makes this qualification beside (d). There is no need to strengthen the main proposition to a saturation condition involving localization.

The subsequent corollaries must be read with their respective closure hypotheses. If $TC$ is retract closed, then $\widetilde{TC}=TC$, so the recognition condition becomes $C\cap W=TC$. The analogous statement for $TF$ needs retract closure of $TF$. If $C$ and $W$ are retract closed, there is an even shorter proof of Corollary~2 than the printed one: $TC\subset C\cap W$ gives
\[
\widetilde{TC}=\operatorname{Retr}(TC)\subset C\cap W,
\]
and condition (a) supplies the reverse inclusion. The printed relation $C\cap W\longleftrightarrow TC$ in this proof is a slip. The lifting relation available here is $C\cap W\longleftrightarrow F$, and the relevant pr\'ecouple is $(C\cap W,F)$.
\end{dcomment}

\begin{dcomment}{development}{p36}{The missing closure hypothesis in the footnote}
Grothendieck's main closure result, Corollary~3, deliberately replaces the four provisional classes by their retract closures. Footnote~13 then proposes a shorter Corollary~4 in which that replacement disappears. As printed, its equivalence with a \emph{closed} Quillen structure is false. The following example isolates the defect while retaining invariance of the classes under isomorphism; it is not an objection based merely on a choice of representatives.

\textbf{A pr\'ecouple whose two classes have not both been closed.}\ours
In the category of sets, with choice, let $W$ be all maps. Let $L$ consist of the injections $i:A\hookrightarrow B$ for which $B\setminus i(A)$ is empty or infinite, and let $R$ consist of all surjections. Set
\[
C=TC=L,\qquad F=TF=R.
\]
Then both required pairs are pr\'ecouples, and conditions (a), (b), and (c) of the printed Corollary~4 hold, but $(W,C,F)$ is not closed.

To verify the lifting condition, consider an injection $i:A\hookrightarrow B$ and a surjection $p:X\twoheadrightarrow Y$. A commutative square specifies the filler on $i(A)$; on each remaining element of $B$, choose an inverse image under $p$ of its prescribed image in $Y$. This proves $L\longleftrightarrow R$. For factorization, take any $f:A\to B$ and form
\[
A\ \longrightarrow\ A\amalg(B\times\mathbb N)
  \ \xrightarrow{\ [f,\operatorname{pr}_B]\ }\ B.
\]
The first map has empty complement when $B=\varnothing$ and infinite complement otherwise; the second is surjective. Thus it is an $(L,R)$-factorization. Since $W$ contains every map, the equalities $TC=C\cap W$ and $TF=F\cap W$ hold, and the factorization just given proves $W=TF\circ TC$.

On the other hand, $\varnothing\to\{0\}$ is a retract of $\varnothing\to\mathbb N$: use the inclusion selecting $0$ and the unique map $\mathbb N\to\{0\}$. The larger map belongs to $L$, whereas the smaller one does not. Therefore $L$ is not retract closed and cannot equal ${}^{\square}R$. This proves the failure of (d).

The appropriate repair of the footnote is to require the two input pairs to be \emph{couples}, or equivalently to require all four classes to be retract closed. Proposition~4.3 then gives the equivalence of (a), (b), and (c), and the resulting two weak factorization systems give (d), subject to the required existence of limits and colimits. Another valid conclusion under the weaker hypotheses is the one actually stated in Corollary~3: take retract closures first. The general modern formulation, including automatic retract closure of $W$ once the two weak factorization systems exist, is in \citet[Definition~11.3.1 and Lemma~11.3.3]{Riehl2014}.

We checked those general criteria and the recognition theorems cited below. We did not find this particular isomorphism invariant example in those sources; the claim here is the explicit counterexample and its verification, not a claim of priority for the retract obstruction.
\end{dcomment}

\begin{dcomment}{exposition}{p39}{Why preservation under pushout is the right condition}
The second factorization cannot be obtained merely by declaring $TC=C\cap W$. Its left class must survive the operations used in the small object construction. Composition and retracts are already available, but an arbitrary weak equivalence need not remain a weak equivalence after pushout. Grothendieck's choice $C_0\subset\Cof_W$ makes the missing stability a property of the cofibrations along which one pushes out.

With the definition recalled in Chapter~XIII, manuscript page~8, a map $i:A\to B$ lies in $\Cof_W$ when its pushout functor
\[
i_!:A\backslash\M\longrightarrow B\backslash\M,
\qquad (A\to X)\longmapsto(B\to B\amalg_A X),
\]
preserves weak equivalences. All pushouts mentioned here are assumed to exist. This is often called the condition of being an $h$-cofibration relative to $W$. The property automatically passes to cobase changes, by the pasting law for pushout squares. It is defined from $W$, before a model structure has been chosen.

The distinction between preserving weak equivalences under pushout \emph{along} $i$ and having every pushout \emph{of} $i$ belong to $W$ is essential. The latter condition defines $W^{\univ}$. Two-out-of-three gives $W^{\univ}\subset\Cof_W\cap W$: in the square comparing the cobase changes of $i$ to $X$ and to $Y$, the two cobase changes and the given map $X\to Y$ are weak equivalences. The reverse inclusion needs a factorization hypothesis. Grothendieck had explicitly used that hypothesis in Chapter~XIII, Corollary~1 and Proposition~4, manuscript pages~10--11; it has been dropped from Corollary~4.4.3 here. The next block records the resulting counterexample and restores the hypothesis exactly where it is needed.

The filtered-colimit assertion of Proposition~4.4.2 itself is sound and has a short explicit proof. Let $i_t:A_t\to B_t$ be a filtered diagram in $\Cof_W$, with colimit $i:A\to B$. Given a weak equivalence $u:X\to Y$ under $A$, restrict it under each $A_t$. Each map
\[
B_t\amalg_{A_t}X\longrightarrow B_t\amalg_{A_t}Y
\]
lies in $W$. Their filtered colimit is $B\amalg_A X\to B\amalg_A Y$, since colimits commute with colimits. Closure of $W$ under filtered colimits proves the required preservation. It also gives retract closure of $W$: a retract is the colimit of the countable sequence which repeats the associated idempotent. Thus a retract of a map in $\Cof_W$ again preserves $W$, because its pushout comparisons are retracts of the corresponding comparisons for the larger map. Together with pushout and composition, these facts prove the asserted stability under continuous transfinite composition by induction on the ordinal.
\end{dcomment}

\begin{dcomment}{development}{p40}{A missing hypothesis in the couniversal-equivalence assertion}
The intended application will have a useful factorization available before the second small object argument starts. Remembering that factorization repairs Corollary~4.4.3 and preserves the eventual recognition theorem. Forgetting it makes the intermediate assertion false even in a finite lattice.

\textbf{A finite counterexample.}\ours
Let $\M$ be the diamond lattice, viewed as a category, with $0<a<1$ and $0<b<1$, and with $a,b$ incomparable. Let $W$ consist of the identity maps and $i:0\to a$. Then $W$ satisfies two-out-of-three and is closed under all filtered colimits; all objects of $\M$ are finitely presentable. Nevertheless,
\[
i\in\Cof_W\cap W,\qquad i\notin W^{\univ}.
\]
Indeed, a filtered diagram in a finite poset has a greatest element in its image, and that element is its colimit. The same observation applies to the finite arrow poset, proving closure of $W$ under filtered colimits. It also shows directly that each representable preserves filtered colimits. To check two-out-of-three, the only nonidentity composable pairs are the two chains from $0$ to $1$ through $a$ and $b$; neither chain has two of its three maps in $W$.

Pushout along $i$ is the operation $x\mapsto a\vee x$. Among the arrows under $0$, the only nonidentity weak equivalence is $0\to a$, and it is sent to $\id_a$. Hence $i\in\Cof_W$. But its pushout along $0\to b$ is $b\to1$, which does not belong to $W$. This verifies every claimed property and disproves the equality in Corollary~4.4.3 under its printed hypotheses. If $C_0=\{i\}$, then $C=\widetilde{C_0}$ contains $i$ and its pushout $b\to1$. Thus $TC=C\cap W$ contains $i$ but not that pushout, disproving also the asserted pushout stability of $TC$ in Corollary~4.4.4. The other identities involving $W^{\univ}$ in that corollary require the same repair.

\textbf{The factorization that repairs the argument.}
Assume that every map factors as a member of $\Cof_W$ followed by a member of $W$. Then $\Cof_W\cap W=W^{\univ}$. This is the dual form of \citet[Chapter~XIII, Proposition~4, ms.~pp.~10--11]{GrothendieckDerivateursXIII}, and we give the complete argument in the notation used here.

Take $i:A\to B$ in $\Cof_W\cap W$ and any $u:A\to X$. Factor $u=qj$ with $j:A\to Z$ in $\Cof_W$ and $q:Z\to X$ in $W$. Form $P=B\amalg_A Z$ and $Q=B\amalg_A X$. The map $k:Z\to P$ is weak because it is the pushout of the weak equivalence $i$ along the $W$-cofibration $j$. It is also a $W$-cofibration, being a pushout of $i$. Therefore the induced map $q':P\to Q$, the pushout of $q$ along $k$, is weak. The equation $q'k=i'q$, where $i':X\to Q$ is the desired cobase change, and two-out-of-three imply $i'\in W$. The reverse inclusion was proved in the preceding block.

Proposition~4.4.6 supplies exactly this hypothesis: its first factorization has left map in $C\subset\Cof_W$ and right map in $TF\subset W$. Once that first factorization is constructed, the proof above establishes the pushout stability of $TC$, and all subsequent uses of Corollary~4.4.4 are justified. The false intermediate generalization therefore does not invalidate the final recognition proposition. The class $W^{\univ}$ itself remains closed under filtered colimits and the stated left stability operations even without the extra factorization.

Our literature check covered the earlier manuscript argument and the modern comparison in \citet[Lemma~1.10]{BataninBerger2017}, which distinguishes couniversal weak equivalences from trivial $h$-cofibrations and gives equality in left proper model categories. We did not find this finite lattice example in those sources.
\end{dcomment}

\begin{dcomment}{examples}{p41:La première assertion est, on l'a vu tantôt}{An inclusion which cannot be an equality}
Grothendieck has established enough closure to apply the small object theorem to a chosen set $C_0\subset\Cof_W$. The conclusion is that the class it generates stays inside $\Cof_W$. The display in Corollary~4.4.4 prints $C=\Cof_W$, although the proof on the following manuscript page proves $C\subset\Cof_W$. This is a slip in the edition, not in the transcription.

A small example shows why the stronger assertion cannot follow. Take $\M=\mathrm{Set}$, let $W$ consist of the bijections, and take $C_0=\{\id_{\{0\}}\}$. Every map belongs to $\Cof_W$, since every pushout functor preserves isomorphisms. Every map also has the right lifting property against $\id_{\{0\}}$, so
\[
C_0^{\square}=\operatorname{Mor}(\mathrm{Set}),\qquad
{}^{\square}(C_0^{\square})=\operatorname{Iso}(\mathrm{Set}).
\]
For the second equality, a map which lifts against every map lifts against itself in the square with identity horizontal maps, and the resulting filler is a two-sided inverse. Thus $C$ consists of the bijections, whereas $\Cof_W$ contains, for example, $\varnothing\to\{0\}$. All objects of $\mathrm{Set}$ are presentable for some regular cardinal, and the required colimits exist, so the example meets the hypotheses of the corollary.

Under the extra factorization hypothesis explained in the preceding block, the rest of the corollary has its intended force:
\[
TC=C\cap W=C\cap W^{\univ}.
\]
Both factors in the last intersection have the required stability under pushout, retracts, and continuous transfinite composition. Hence so does $TC$. This establishes the necessary closure of the proposed trivial cofibrations; the existence of a small set generating them remains the next problem.
\end{dcomment}

\begin{dcomment}{exposition}{p42:Il reste à prouver l'équivalence}{The two identities already available before the second factorization}
Grothendieck's two footnotes on this page separate a formal recognition problem from a problem of size. Put $C=\widetilde{C_0}$, $TF=C^{\square}\subset W$, $TC=C\cap W$, and $F=TC^{\square}$. The first weak factorization system $(C,TF)$ already gives
\[
W=TF\circ TC,\qquad TF=F\cap W,
\]
without assuming a factorization with left map in $TC$ for every arrow. The missing quantifier concerns arbitrary arrows, not only weak equivalences.

Here are complete proofs. Given $w\in W$, factor it as $w=pi$ using $(C,TF)$. Since $p\in TF\subset W$, two-out-of-three gives $i\in C\cap W=TC$. Conversely, a composite of a member of $TC$ and a member of $TF$ belongs to $W$. For the second identity, $TF\subset F\cap W$ follows from $TC\subset C$. If $f\in F\cap W$, factor $f=pi$ with $i\in C$ and $p\in TF$. Again $i\in TC$. The square with left arrow $i$, right arrow $f$, top arrow the identity of the source of $f$, and bottom arrow $p$ has a diagonal $s$ satisfying $si=\id$ and $fs=p$. Consequently $f$ is a retract of $p$ in the arrow category. The right lifting class $TF$ is retract closed, so $f\in TF$.

For a candidate set $TC_0\subset TC$, the notation left implicit in the statement of Proposition~4.4.6 is $F_0=TC_0^{\square}$. Smallness produces a weak factorization system $(\widetilde{TC_0},F_0)$, but only the inclusions
\[
\widetilde{TC_0}\subset TC,\qquad F\subset F_0
\]
are automatic. Proposition~4.4.6 says that the candidate succeeds precisely when $F_0\cap W=TF$, equivalently when every weak equivalence has the more restrictive factorization with left map in $\widetilde{TC_0}$. If $F_0=F$, then $\widetilde{TC_0}={}^{\square}F=\widetilde{TC}$; combined with $\widetilde{TC_0}\subset TC\subset\widetilde{TC}$, this gives equality immediately. Thus the final implication does not assume the desired second factorization in disguise.

To see the distinction concretely, take $\M=\mathrm{Set}$ and $W$ all maps. Let $C_0$ contain $\varnothing\to\{0\}$ and the fold $\{0,1\}\to\{0\}$. Right lifting against the first map is surjectivity, and right lifting against the second is injectivity. Hence $TF$ consists of bijections and $C$ consists of all maps. Then $TC$ is all maps and $F$ consists of bijections. If the candidate $TC_0$ contains only $\varnothing\to\{0\}$, its right class $F_0$ consists of surjections. The projection $\{0,1\}\to\{0\}$ belongs to $F_0\cap W$ but not to $TF$. Both identities of the preceding paragraph still hold; the candidate has simply omitted a necessary generating map.
\end{dcomment}

\begin{dcomment}{later}{p43:J'ai passé des jours à essayer}{The generating set supplied by Smith's theorem}
The question at the end of Section~4 became a standard recognition theorem for combinatorial model categories. In the locally presentable setting, Smith's theorem, published with proof by \citet[Theorem~1.7]{Beke2000}, gives the missing set of generating trivial cofibrations. Its hypotheses are precisely organized around the difficulty seen here: $W$ has two-out-of-three and is retract closed, $C_0^{\square}\subset W$, $C\cap W$ is stable under pushout and transfinite composition, and $W$ satisfies a solution set condition at each map of $C_0$. Accessibility of the full subcategory $W\subset\M^{[1]}$ supplies the last condition by \citet[Proposition~1.15]{Beke2000}.

The solution set condition says that, for each $i\in C_0$, a set of maps $i\to w_s$ with $w_s\in W$ is weakly initial among all maps from $i$ to members of $W$ in the arrow category. This is the categorical form of the bounded choices Grothendieck is seeking. Smith's argument converts these choices into a set $J\subset C\cap W$ with $\operatorname{cof}(J)=C\cap W$; Beke's Lemmas~1.8 and~1.9 supply the construction and proof.

This answers the existence question under the stated modern hypotheses. A cardinal filtration on its own should not silently be substituted for local presentability, and the claim does not cover every category allowed by the more flexible assumptions on the preceding pages. Nor does the theorem require $C$ itself to be stable under all filtered colimits, the additional route Grothendieck begins to explore in Section~5.
\end{dcomment}

\begin{dcomment}{status}{p43:J'ai passé des jours à essayer}{The form closest to the proposed localizer construction}
Grothendieck's use of $\Cof_W$ also survives in a recognition theorem whose hypotheses closely follow these pages. In the updated online text of \citet[Proposition~A.2.6.15]{Lurie2009HTT}, a locally presentable category, a perfect class $W$, and a set of generating cofibrations satisfying the pushout preservation condition and $C_0^{\square}\subset W$ produce a left proper combinatorial model structure. Here ``perfect'' means that $W$ contains the isomorphisms, satisfies two-out-of-three, is closed under filtered colimits, and is generated under those colimits by a set. The proof first establishes the required stability of $C\cap W$, then applies the small-generation result of Lemma~A.2.6.9 and the general recognition Proposition~A.2.6.10. These are the numbers in the online revision consulted, not the earlier printed numbering.

The later theorem makes precise both parts of the intended conclusion: the second factorization exists, and the chosen cofibrations preserve weak equivalences under pushout. The latter is left properness. It remains separate from the formal two-factorization axioms studied at the beginning of Section~4; an arbitrary model structure need not have it.
\end{dcomment}


\begin{dcomment}{later}{p43:J'ai passé des jours à essayer}{The passage from models to diagram homotopy theories}
The point of obtaining a model structure is that it can support coherent operations on diagrams. In the combinatorial setting, the later passage is precise: \citet[Theorem~1.28 and Proposition~1.30]{Groth2013} constructs projective and injective model structures on diagram categories and proves that $J\mapsto\operatorname{Ho}(\mathcal M^J)$ is a strong derivator. The two diagram model structures supply the appropriate derived Kan extensions.

The existence of these diagram model structures is a convenient sufficient route. The broader construction of homotopy right Kan extensions for model categories with small limits is treated by \citet[Theorem~6.11]{Cisinski2003}, in that paper's left-derivator convention. This later theorem explains why the goal is larger than a particular generating-set argument. Chapter~XV builds and controls models; the derivator records their homotopy theories of diagrams and the functors between them. Neither the cardinal estimates nor the two lifting systems alone constitute the passage between the two.
\end{dcomment}

% Section 5 commentary; GPT-6 Astra (OpenAI).

\begin{dcomment}{exposition}{p44}{Images and the meaning of a small subobject}
Grothendieck now wants the size of an object to be visible inside the object itself. A presentation by a filtered diagram can repeat objects, identify elements, and carry many choices. The ordered collection of its small subobjects is a more geometric presentation: its transition maps are inclusions, and its small directed suprema can be calculated in the ambient category. This change of presentation is the preparation for recognizing a class of objects from its bounded part.

We keep the three conditions on the filtration separate. Essential smallness says that a bounded collection has a set of isomorphism classes; generation says that every object has an appropriately directed presentation by bounded objects; preservation of such colimits by a represented functor is an additional condition on the representing object. In the terminology of these pages, ``grand devant $\pi$'' means that every subset of cardinal at most $\pi$ has an upper bound. For infinite $\pi$, this is the usual $\pi^+$-directedness. A $\pi$-accessible object is consequently a $\pi^+$-presentable object in current terminology. It is not an object with exactly $\pi$ elements, and the bound is not assumed regular.

The set-sized assertion in Lemma~5.1 also has content. Let $G$ run through a small set of representatives of the base level. A subobject $U\hookrightarrow X$ determines the subsets of $\Hom(G,X)$ consisting of the maps which factor through $U$. These subsets determine $U$: a generating presentation of its domain lifts uniquely to any other subobject having the same subsets, and these unique lifts assemble. Thus there is a set of subobjects of $X$.

The same base level is a conservative generating family. Suppose $f:X\to Y$ induces bijections on maps from its objects. Write $Y=\varinjlim_iG_i$ using the filtration. Each structural map $G_i\to Y$ lifts uniquely to $X$; uniqueness makes the lifts compatible, so they induce $g:Y\to X$ with $fg=\mathrm{id}_Y$. For a presentation $X=\varinjlim_jH_j$, injectivity of $\Hom(H_j,f)$ gives $gf h_j=h_j$ at every structural map $h_j$, hence $gf=\mathrm{id}_X$. No accessibility rank is needed. The dual of \citep[Expos\'e~I, 8.12.7]{SGA4I1972} now supplies completeness: the functor of cones $Z\mapsto\varprojlim_j\Hom(Z,X_j)$ takes colimits to limits and is represented by the desired limit. This argument uses conservative generation, without requiring the stronger density assertion in the source.

Once intersections exist, the image used here is the intersection of \emph{all} subobjects through which the map factors. If $Z\to\operatorname{Im}(f)$ were not epic, two different maps out of its codomain would agree after composition; their equalizer would give a smaller subobject through which $f$ factors. This proves Corollary~5.2. In fact the first map in this factorization is extremal epic: any factorization through a monomorphism would contradict the minimality of the image unless that monomorphism were invertible.

In Lemma~5.3 the adjective \emph{strict} has the categorical meaning of SGA~4, not merely ``different from the whole object''; the latter requirement is written separately. A strict monomorphism which is also epic is invertible. Conversely, failure to be epic is detected by a proper equalizer, which is a strict monomorphism. This proves the stated criterion, after replacing the accidental name $g$ by $f$. The distinction matters, for example, for the epimorphism $\mathbb Z\to\mathbb Q$ in commutative rings: it factors through a proper monomorphism, namely itself, but not through a proper strict monomorphism. The terminology is fixed in \citep[Expos\'e~I, 10.2--10.5]{SGA4I1972}.
\end{dcomment}

\begin{dcomment}{development}{p46}{Quotient bounds do not bound subobjects}
The two image conditions serve different purposes. The second controls the image of a bounded source; the first controls an arbitrary subobject of a bounded target. The suggestion in footnote~17 that the second should imply the first is false even in a module category with exact filtered colimits. Giving the example explicitly explains why the two conditions must be distinguished, even in arguments where only one of them is used.

\textbf{A filtration separating $(d_1)$ and $(d_2)$.}\ours
Fix an infinite cardinal $\pi$, a field $k$, and a set $S$ of cardinal greater than $\pi$. Put
\[
 R=k[x_s\mid s\in S],\qquad
 \Filt^{\kappa}(R\text{-}\mathrm{Mod})
 =\{N\mid N\text{ is generated by at most }\kappa\text{ elements}\},
 \quad \kappa\geq\aleph_0.
\]
This is a cardinal filtration satisfying $(a)$, $(b)$, $(c)$, $(d_0)$, $(d_2)$ and $(e)$, but $(d_1)$ fails at $\pi$.

\emph{Proof.} Modules generated by at most $\kappa$ elements are quotients of $R^{(\kappa)}$, so they have a set of isomorphism classes. A colimit of a diagram with at most $\kappa$ arrows and $\kappa$-generated values is a quotient of the sum of at most $\kappa$ such modules; it is again $\kappa$-generated. This proves $(d_0)$ and the required closure part of the filtration, and quotient closure proves $(d_2)$. Every module is presentable for some sufficiently large regular cardinal: choose a presentation, and bound both its generators and relations; factorization and eventual equality in a sufficiently directed colimit are then checked on those sets. This proves $(c)$. Filtered colimits of modules are exact, giving $(e)$.

For the quantitative generation clause, let $N$ have a chosen generating set $T$ of size at most $\kappa'$, where $\kappa'\geq\kappa$. The submodules generated by subsets of $T$ of size at most $\kappa$ form a $\kappa^+$-directed diagram with colimit $N$. Its indexing set has cardinal at most $(\kappa')^{\kappa}$. This is precisely the bound required of a cardinal filtration, including when the ring itself is much larger than $\kappa$.

Finally $R$ is cyclic, whereas its ideal $\mathfrak a=(x_s\mid s\in S)$ is not $\pi$-generated. Indeed $R/\mathfrak a=k$ and
\[
 \mathfrak a/\mathfrak a^2\cong\bigoplus_{s\in S}k\,\overline{x_s}.
\]
Any generating set of $\mathfrak a$ maps to a spanning set of this $k$-vector space and therefore has cardinal at least $|S|$. Thus $\mathfrak a\hookrightarrow R$ disproves $(d_1)$.\hfill$\square$

The example does not challenge the propositions which follow: they assume $(d)$, hence both image conditions. It also explains why increasing the cardinal can remove the difficulty in familiar concrete categories. If $\kappa\geq|R|$, every submodule of a $\kappa$-generated module has at most $\kappa$ underlying elements and is itself $\kappa$-generated. We checked the filtration definition and the relevant size arguments in \citep[Expos\'e~I, 9.12--9.16]{SGA4I1972}; the particular counterexample above is our explicit verification of the remark's failure, with no claim of priority.
\end{dcomment}

\begin{dcomment}{exposition}{p49}{The canonical directed union}
The gain in Proposition~5.3 is a canonical presentation whose indexing order has actual suprema. Write $\operatorname{Sub}_{\pi}(X)$ for Grothendieck's $\Ss_{\pi}(X)$. Given at most $\pi$ members $U_a\hookrightarrow X$, form their coproduct and then its image in $X$. The coproduct belongs to $\Filt^\pi$ by $(d_0)$, and its image belongs to $\Filt^\pi$ by $(d_2)$. It is their least upper bound in the full subobject lattice. These are the two assumptions used in the proof; the printed references to $(c)$ and $(d_1)$ on manuscript page~47 are misplaced.

Here is the colimit argument in a form that also explains the role of epimorphisms. Choose a generating presentation $X=\varinjlim_iX_i$ by $\Filt^\pi$-objects, and factor each map to $X$ as
\[
 X_i\twoheadrightarrow U_i\hookrightarrow X.
\]
A compatible family $U_i\to T$ gives, by composition, a compatible family $X_i\to T$ and hence a unique map $X\to T$. Its restriction to $U_i$ is the given map because $X_i\to U_i$ is epic. Therefore $\varinjlim_iU_i=X$. Passing to the image of the indexing order in $\operatorname{Sub}_{\pi}(X)$ does not change this colimit: above any chosen $U_i$, the original directed order supplies a common upper bound for every finite family of representatives.

There is consequently a section of the canonical map
\[
 q:\varinjlim_{U\in\operatorname{Sub}_{\pi}(X)}U\longrightarrow X.
\]
Under $(e)$ this map is monic, being the filtered colimit of the monomorphisms $U\to X$; a monomorphism with a section is invertible. Alternatively, if every $\Filt^\pi$-object is $\pi$-accessible, any small subobject $V\hookrightarrow X$ factors through some $U_i$, and the factorization is a map of subobjects. The $U_i$ are then cofinal in $\operatorname{Sub}_{\pi}(X)$, proving the same result without $(e)$.

For the stronger continuity assertion of Scholie~5.5, $(e)$ is used again. If $K\subset\operatorname{Sub}_{\pi}(X)$ is directed and $|K|\leq\pi$, its colimit belongs to $\Filt^\pi$ by $(d_0)$, and its map to $X$ is monic by $(e)$. It is the supremum of $K$. Thus the canonical diagram carries small directed suprema to colimits. This is the precise content of a \emph{perfectly $\pi$-adapted} representation. It is stronger than the mere requirement that its index be $\pi^+$-directed.

In a presheaf category on a small category $A$, this construction has a concrete form once $\pi$ bounds $|\operatorname{Fl}(A)|$. Use the filtration by presheaves whose total set of elements has cardinal at most $\pi$. Closing at most $\pi$ chosen elements under all restriction maps still gives at most $\pi$ elements. The poset of such subpresheaves is $\pi^+$-directed, its small directed suprema are unions, and their union is the whole presheaf. This familiar example is the geometric picture underlying the abstract image argument.
\end{dcomment}

\begin{dcomment}{exposition}{p51}{Accessibility tests every adapted presentation}
The word ``accessible'' now applies to a \emph{class of objects}. Its meaning must be recovered from the definition to which Grothendieck refers, since it is more specific than simply saying that the corresponding full subcategory is accessible in the modern sense. The relevant condition $b_\pi$ in \citep[Chapter~XIII, manuscript pp.~88--89]{GrothendieckDerivateursXIII} is this: for every perfectly $\pi$-adapted presentation
\[
 X=\varinjlim_{i\in I}X_i,\qquad X_i\in\Filt^\pi,
\]
one has $X\in M$ if and only if $\{i\mid X_i\in M\}$ is cofinal in $I$. Thus the test concerns all such presentations, including presentations whose maps identify elements.

Formula~(5.1) applies this test to the canonical subobject presentation. It records a necessary condition, and gives a way to reconstruct a candidate class from any chosen bounded part $M^\pi$. On bounded objects the recovery formula is automatic: $X$ itself is the largest element of $\operatorname{Sub}_{\pi}(X)$, so a cofinal subset must contain it. This explains (5.2), but supplies no invariance under changing the presentation.

Condition~(5.3) asks for that invariance among the continuous cofinal subdiagrams of the canonical diagram. The closure requirement on $J$ says that it contains the supremum of every directed subset of itself of cardinal at most $\pi$. It makes the induced diagram perfectly adapted. It does not turn a nonmonic presentation into a subobject presentation, and this is where the first version of Proposition~5.6 overreaches. The source itself identifies the obstruction and inserts the correction in footnote~21.

The repair on manuscript page~55 asks that $M^\pi$ be closed under filtered colimits of diagrams of cardinal at most $\pi$. Then the good indices in a perfectly adapted diagram are themselves closed under the same directed suprema. They form another perfectly adapted presentation whenever they are cofinal. Lemma~5.7 can therefore be applied to that presentation, reducing its colimit to a continuous cofinal family of monomorphic stages. This is why the extra closure condition repairs the converse: it allows the comparison lemma to be used a second time, on the good part of the diagram.
\end{dcomment}

\begin{dcomment}{development}{p55}{A two-point model of the retract obstruction}
Grothendieck's proposed counterexample asks for an object which has no proper monic self-map, but does have a nontrivial retract. The walking arrow provides such an object in an elementary ambient category, and also makes the failure visible in a perfectly adapted presentation. This gives a concrete completion of the mechanism proposed in (5.4).

\textbf{A class satisfying (5.3) but failing $b_\pi$.}\ours
Let $\mathcal M=\mathrm{Pos}$ be the category of small partially ordered sets and monotone maps. For an infinite cardinal $\pi$, use the filtration by underlying cardinality at most $\pi$, and let $M$ consist of the objects isomorphic to the two-element chain $P=\{0<1\}$. Then $M$ satisfies (5.3) but is not $\pi$-accessible.

\emph{Proof.} Monomorphisms in $\mathrm{Pos}$ are the injective monotone maps, so every monic endomorphism of $P$ is the identity. The constant map $P\to P$ with value $0$ is nevertheless an idempotent whose image is the one-point poset $1$. The ambient hypotheses are satisfied: a colimit is obtained from a disjoint union by generating the order relations and quotienting by the resulting equivalence relation, so a diagram of cardinal at most $\pi$ with $\pi$-small values has a $\pi$-small colimit. Subobjects and image quotients have the same cardinal bound. Filtered colimits have underlying set the filtered colimit of the underlying sets, and an order relation between two representatives holds exactly when it holds at some later common stage. This description proves $(e)$. A poset of cardinal at most $\pi$ has at most $\pi$ order relations; choosing representatives for its elements and making those relations hold at one common stage proves its $\pi$-accessibility. The eventual-equality part follows by the same argument on the at most $\pi$ elements.

To check (5.3), first suppose $X\cong P$. Every cofinal subset of $\operatorname{Sub}_{\pi}(X)$ contains its largest member $X$, so its good members are cofinal. Conversely, suppose $J^M$ is cofinal in a continuous cofinal subset $J$ of $\operatorname{Sub}_{\pi}(X)$. Then $X$ is the filtered colimit of the subobjects indexed by $J^M$. Each has type $P$, and every transition between them is an isomorphism. The index is nonempty and directed, so its colimit is again $P$. This establishes both directions of (5.3).

For a direct failure of the actual definition $b_\pi$, put $\lambda=\pi^+$ and index a diagram by the ordinals $\alpha<\lambda$. Take $X_\alpha=P$ at $0$ and at every successor ordinal, and $X_\alpha=1$ at every nonzero limit ordinal. Give each $X_\alpha$ its distinguished least element. For $\alpha<\beta$, let the transition map be the constant map to that element of $X_\beta$. These maps compose as required. The colimit is $1$, since every element is identified with a distinguished element at a later stage.

This diagram is perfectly $\pi$-adapted. The ordinal $\lambda$ is $\pi^+$-directed. Every nonempty directed subset of cardinal at most $\pi$ has a supremum below $\lambda$; if it has a maximum, the colimit is its final stage, and otherwise its supremum is a nonzero limit ordinal and both the colimit and that stage are $1$. Yet the good indices, namely $0$ and the successor ordinals, are cofinal, while the colimit is not in $M$. Thus $b_\pi$ fails.\hfill$\square$

The same phenomenon can be described as splitting an idempotent by a long filtered diagram. The continuous version above shows why replacing arbitrary presentations by subobject presentations loses precisely the needed information. We checked the definition in Chapter~XIII, manuscript pp.~85--97, and the treatment of idempotents in \citep[Expos\'e~I, 8.6.5 and 10.6]{SGA4I1972}. We did not find this specified two-point diagram in those sources; it is an elaboration of the manuscript's counterexample, not a claim of priority.
\end{dcomment}

\begin{dcomment}{exposition}{p57}{Making the image factorization a transition}
Lemma~5.7 is the comparison that makes the repaired recognition theorem work. Its diagonal construction is sound once one supplies an eventual-equality step. That step uses the accessibility of $X_i$ as well as the accessibility of its image; factorization of the image alone does not identify the resulting composite with a transition of the original diagram.

Write $q_i:X_i\twoheadrightarrow U_i$ for the image factorization of the structural map to $X$. Given $i_0$, accessibility of $U_{i_0}\in\Filt^\pi$ gives $j\geq i_0$ and a factorization
\[
 U_{i_0}\xrightarrow{a}X_j\longrightarrow X.
\]
The two maps $a q_{i_0}$ and $X(i_0\leq j)$ from $X_{i_0}$ to $X_j$ have the same composite to $X$. Accessibility of $X_{i_0}$ therefore gives a further stage $i_1\geq j$ at which these maps agree. After replacing $a$ by its composite to $X_{i_1}$, the upper composite in Grothendieck's zigzag really is the transition $X_{i_0}\to X_{i_1}$. Its lower composite is the inclusion $U_{i_0}\hookrightarrow U_{i_1}$, since both maps are over $X$ and the latter image maps monomorphically to $X$.

Repeating gives an interleaving of the chains $(X_{i_n})$ and $(U_{i_n})$. A cocone on either chain extends uniquely to a cocone on the other, so their colimits agree. Let $i_\omega=\sup_n i_n$. Perfect adaptation identifies $X_{i_\omega}$ with that colimit, while $(e)$ makes the map $\varinjlim_nU_{i_n}\to X$ monic. Thus $i_\omega$ is a monomorphic stage above $i_0$. The same use of $(e)$ and continuity proves that the monomorphic stages are closed under directed suprema of cardinal at most $\pi$. The choices may be non-strictly increasing; the proof does not require that the index have no maximal element.

There is a second detail in applying this lemma on manuscript page~53. The image of a continuous cofinal map of directed posets need not itself be closed under directed suprema, as the example below shows. Consequently, the image $J$ of the monomorphic stage indices cannot simply be inserted into (5.3) without checking its closure requirement.

The correction already adopted on manuscript page~55 supplies the repair. If $X\in M$, the good small subobjects form a cofinal subset closed under small directed unions, because $M^\pi$ has that closure. Starting with any monomorphic stage, choose a good small subobject containing it, then a later monomorphic stage containing that subobject, and continue alternately. At the countable supremum the stage and the union of the good subobjects agree, so the resulting stage is good. Hence good stages are cofinal. In the converse direction, if good stages are cofinal, their closure under small directed suprema makes their restricted diagram perfectly adapted. Applying Lemma~5.7 to it produces cofinally many good monomorphic stages; their subobjects are cofinal in $\operatorname{Sub}_{\pi}(X)$, so (5.1) gives $X\in M$. This proves the corrected assertion without imposing a closure property on the raw image of an indexing map.
\end{dcomment}

\begin{dcomment}{development}{p57}{A continuous image need not be closed}
The distinction just used is easy to miss because the indexing map preserves every supremum which can be formed upstairs. The difficulty is that a directed family of image values need not have coherently ordered representatives. The following example isolates that point in sets.

\textbf{A perfectly adapted diagram with nonclosed image.}\ours
Let $I$ be the countable subsets of the successor ordinals below $\omega_1$, ordered by inclusion. Choose $e\notin\omega_1$, put $X=\omega_1\cup\{e\}$, and define subobjects of $X$ by
\[
 X_S=\bigcup_{\alpha\in S}\alpha
       \ \cup\ \{e\mid |S|\geq2\}.
\]
With inclusion maps this is a perfectly $\aleph_0$-adapted presentation of $X$, and its image is cofinal among the countable subobjects of $X$, but that image is not closed under countable directed suprema.

\emph{Proof.} The union of countably many countable subsets is countable, so $I$ is $\aleph_1$-directed, and countable directed suprema in $I$ are unions. If $\mathcal K\subset I$ is directed, its union contains two distinct ordinals exactly when some member of $\mathcal K$ contains both: direct the two witnesses to a common member. It follows that
\[
 X_{\bigcup\mathcal K}=\bigcup_{S\in\mathcal K}X_S.
\]
Every $X_S$ is countable, and the union of the whole diagram is $X$, proving perfect adaptation. Any countable subset of $X$ is contained in $X_S$ for a two-element $S$ containing a successor ordinal above its ordinal elements; this proves cofinality. Finally every positive finite ordinal $n$ is $X_{\{n\}}$. Their directed supremum is $\omega$ without $e$. No $S$ with two or more elements gives this value, because its value contains $e$; a singleton gives a successor ordinal, and the empty set gives the empty set. Thus $\omega$ is not in the image.\hfill$\square$

The literature check here was bounded to the adapted-presentation discussion in Chapter~XIII, manuscript pp.~85--97, and SGA~4, Expos\'e~I, 9.1--9.3 and 9.12--9.16. This explicit example supplies a missing verification in the present argument; we make no assertion that the observation is new.
\end{dcomment}

\begin{dcomment}{exposition}{p60}{Closed cofinal families and the four descriptions}
The alternating construction in Lemma~5.8 is a small version of the closed-unbounded argument familiar in set theory. Its purpose here is representation independence. Two canonical-looking systems of small subobjects need not meet cofinally merely because each is cofinal. Closure under countable directed suprema is what allows an alternating sequence to finish in both systems.

For example, the even and odd natural numbers are cofinal in $\mathbb N$ and disjoint. That order lacks the required countable suprema. In a $\pi$-adapted order, with $\pi$ infinite, start above any prescribed element and alternate choices in two closed cofinal subsets. The supremum of the even subsequence equals that of the odd subsequence, and belongs to both. The intersection is therefore cofinal; closure of the intersection is immediate. This gives exactly the assertion used in the manuscript, without any assumption that $\pi$ itself is regular.

We can now organize the corrected Proposition~5.7. Let $S\subset\Filt^\pi$ be invariant under isomorphism and closed under filtered colimits with at most $\pi$ arrows. Define $M_S$ by cofinality of its $S$-subobjects in $\operatorname{Sub}_{\pi}(X)$. Under $(e)$ these good subobjects are closed under small directed suprema. Their intersection with any other closed cofinal family is cofinal by the preceding argument. Thus $M_S$ satisfies (5.3), and the corrected Proposition~5.6 identifies it as a $\pi$-accessible class when the $\Filt^\pi$-objects are $\pi$-accessible. Proposition~5.9, in this same range of cardinals, supplies closure under all filtered colimits.

Conversely, a $\pi$-accessible class closed under all filtered colimits has bounded part $S=M\cap\Filt^\pi$ with the required small-colimit closure, by $(d_0)$, and (5.1) reconstructs $M$. These two operations are inverse because on a bounded object cofinality is equivalent to containing the largest subobject, namely the object itself. This is Corollary~5.8.

The description by an arbitrary generating set $M_0\subset\Filt^\pi$ follows by first closing $M_0$ under the filtered colimits of cardinal at most $\pi$ which exist inside the bounded level, and under isomorphism. Call the resulting class $S$. It remains essentially small because it is contained in $\Filt^\pi$. The class $M_S$ is filtered-colimit closed and contains $M_0$. Any filtered-colimit closed class containing $M_0$ contains $S$, and then contains $M_S$, since every object of $M_S$ is a filtered colimit of its good small subobjects. Hence $M_S$ is precisely the filtered-colimit closure of $M_0$.

The printed condition $(c')$, and the matching parenthesis in Corollary~5.8, omit the words saying that the smallest closed class must \emph{contain} the proposed generators. Those words are indispensable: without them the smallest class closed under nonempty filtered colimits is the empty class. The intended statement and its proof are unambiguous from the surrounding discussion. The stronger question in footnote~22 concerns retaining the same rank without the special assumption on the bounded objects; it is not resolved merely by the proof at the admissible ranks used in the proposition.
\end{dcomment}

\begin{dcomment}{exposition}{p64}{The slice category that must be integrated}
The final argument returns from subobjects to maps, because maps assemble functorially when the target moves. Its fibre over $Y\in M$ is $M^\pi/Y$, the category of maps $A\to Y$ with $A\in M^\pi$. The poset naturally included in this fibre is
\[
 \operatorname{Sub}_{\pi}^{M}(Y)
 =\{U\hookrightarrow Y\mid U\in M\cap\Filt^\pi\}.
\]
The unrestricted $\Ss_\pi(Y)$ printed in Lemma~5.10 and its first display must be restricted in this way. An arbitrary small subobject of an object of $M$ need not itself belong to $M$, so the printed unrestricted functor is not defined in general.

Under the image and continuity hypotheses used earlier in this section, the corrected assertion follows directly from (5.1). Given $A\to Y$, its image is small by $(d_2)$, and is contained in a good small subobject by cofinality. Hence every object of $M^\pi/Y$ maps to an object of $\operatorname{Sub}_{\pi}^{M}(Y)$. There is at most one map over $Y$ into any fixed subobject. The good subobjects form a directed poset, and their inclusion in the slice is final: the relevant comma category is nonempty, and two of its objects have a common upper bound. This also proves that $M^\pi/Y$ is filtered. Its source diagram has colimit $Y$, because the final subobject diagram does.

For a filtered diagram $i\mapsto Y_i$ in $M$, the total category $\mathcal E$ has objects $(i,A,a)$, where $A\in M^\pi$ and $a:A\to Y_i$. A morphism to $(j,B,b)$ is a pair $(u,v)$ with $u:i\to j$, $v:A\to B$, and $bv=Y(u)a$. Postcomposition gives an opfibration $\mathcal E\to I$. The adjective ``fibrant'' in Corollary~5.11 has the wrong variance for this construction; no pullback closure of $M^\pi$ is assumed.

The total category is filtered. Two objects can first be moved to a common base index and then given a common successor in that fibre. For two parallel arrows, first equalize their base arrows in $I$, then equalize the resulting two vertical arrows in the target fibre. The same argument shows that $\mathcal E\to I$ is final. More particularly, a cocone on the source objects $A$ of $\mathcal E$ is exactly a compatible family of cocones on the fibres. Taking their colimits gives
\[
 \varinjlim_{(i,A,a)\in\mathcal E}A
 \cong \varinjlim_{i\in I}
          \left(\varinjlim_{(A\to Y_i)\in M^\pi/Y_i}A\right)
 \cong \varinjlim_iY_i.
\]
This is an ordinary colimit calculation; no homotopy-cofinality theorem is being used.

Corollary~5.11 prints $Y_i\in M^\pi$, a stronger hypothesis than the argument needs. The same proof works for $Y_i\in M$, and this is the form required on manuscript page~62. If the base index is a poset, the full subcategory consisting of good subobjects in each fibre is again a poset and is final in $\mathcal E$. Its projection need not remain an opfibration: the image of a good subobject after postcomposition need not be good. The proof requires only finality of this subcategory, which the preceding factorization through good subobjects establishes.
\end{dcomment}

\begin{dcomment}{exposition}{p65}{Completing the filtered-colimit closure argument}
The last reindexing in Proposition~5.9 is the right way to pass from bounded filtered colimits to arbitrary ones. There is, however, a scope issue in the preceding reduction. The proposition drops the image hypotheses $(d_1),(d_2)$ and the exactness hypothesis $(e)$, whereas the proof of Lemma~5.10 as printed uses the earlier subobject presentation. We give a complete proof under the additional assumption that every $\Filt^\pi$-object is $\pi$-accessible. This is the assumption already present in Proposition~5.7 and suffices for every application to the four equivalent descriptions above. The printed general assertion for an arbitrary rank is not certified by the image proof alone.

Put $A=M\cap\Filt^\pi$, and suppose that $A$ is closed under filtered colimits indexed by categories with at most $\pi$ arrows. When forming slices, replace this essentially small class by a small full set of representatives; this does not change their colimits. The necessary direction of Proposition~5.9 is immediate from $(d_0)$, so consider a filtered diagram $Y_i\in M$ with colimit $Y$. For each $Y_i$, choose a perfectly $\pi$-adapted presentation by objects of $\Filt^\pi$. Such presentations exist under $(b),(d_0)$: apply the reindexing by small filtered subcategories given below to any bounded generating presentation. Condition $b_\pi$ makes its $A$-stages cofinal. Thus $Y_i$ is a $\pi^+$-directed colimit of objects of $A$.

Smallness of the domains now supplies the slice argument without images. Any finite diagram in $A/Y_i$ factors through one stage of the chosen $A$-presentation: factor the finitely many structural maps, choose a common stage, and then pass further until its finitely many commutativity equations hold. The same reasoning works for a diagram with at most $\pi$ arrows. Hence $A/Y_i$ is $\pi^+$-filtered. The chosen presentation is final in it, since every map from an $A$-object factors through a stage and two factorizations become equal at a later stage. In particular
\[
 Y_i\cong\varinjlim_{A/Y_i}A.
\]
The opfibration argument in the preceding block then expresses $Y$ as a filtered colimit of $A$-objects, say $Y=\varinjlim_{e\in E}A_e$ with $E$ small and filtered.

To make this presentation perfectly adapted, let $P$ be the poset of filtered subcategories $J\subset E$ with at most $\pi$ arrows. Every set of at most $\pi$ arrows of $E$ is contained in such a subcategory: repeatedly adjoin identities, composites, cocones for finite diagrams and equalizers of parallel pairs. There are at most $\pi$ choices at each step, and countably many repetitions suffice, since $\pi$ is infinite. This also shows that $P$ is $\pi^+$-directed. A directed family of at most $\pi$ members of $P$ has its union as supremum; that union is still filtered and has at most $\pi$ arrows.

Define $A_J=\varinjlim_{e\in J}A_e$. By the assumed bounded closure, $A_J\in A$. Colimits commute with this directed union of indexing categories: a cocone on the union is exactly a compatible family of cocones on the pieces. Consequently
\[
 Y\cong\varinjlim_{J\in P}A_J,
 \qquad
 A_{\bigcup K}\cong\varinjlim_{J\in K}A_J
 \quad(K\subset P\text{ directed},\ |K|\leq\pi).
\]
This is a perfectly $\pi$-adapted presentation with every stage in $M$. Applying $b_\pi$ gives $Y\in M$, which completes the proof in the stated range.

The proof also makes the remaining issue precise. To remove the extra presentability assumption one needs, under the proposition's weaker hypotheses, a replacement for the slice factorization and density argument. The statement on manuscript page~44 about generation by strict epimorphisms would provide such density if established at the required bounded level, but it should not be inferred merely from a choice of colimit presentations. Density and generation under colimits are distinct notions. The earlier discussion already flags this difficulty: Chapter~XIII, manuscript p.~92, questions the asserted strict generation, and footnote~44 on manuscript p.~93 explicitly identifies the tacit hypothesis that the bounded objects are $\pi$-accessible \citep{GrothendieckDerivateursXIII}. We have not settled that unrestricted fixed-rank step here; this is a limit of the present verification, not a claim that it is an open problem in the literature.
\end{dcomment}

\begin{dcomment}{later}{p66}{Filtered-colimit classes and modern accessible categories}
Definition~5.12 isolates the feature needed for the intended application: a class is generated by bounded objects and is closed under \emph{all} filtered colimits. At the good cardinals used in Proposition~5.7, its full subcategory is $\pi^+$-accessible in the current sense, and its inclusion into $\mathcal M$ preserves filtered colimits. Indeed its bounded objects are $\pi^+$-presentable even in the ambient category, every object has a $\pi^+$-directed presentation by them, and the required colimits exist in the full subcategory because they are computed in $\mathcal M$. This is the familiar small-generation description of accessible categories; see \citep[Definition~2.1 and Theorem~2.26]{AdamekRosicky1994}. The manuscript's $b_\pi$ also controls every continuous bounded presentation, so one should not identify its fixed-rank definition with the modern definition without the comparison just proved.

A simple example shows that L-accessibility does not mean local presentability of the full subcategory. In $\mathrm{Set}$, let $M$ be the nonempty sets. Its bounded part consists of the nonempty sets of cardinal at most $\pi$ and is closed under filtered colimits of cardinal at most $\pi$. For a nonempty $X$, every small subset can be enlarged to a nonempty small subset; for $X=\varnothing$ there is no good subobject at all. Thus the small-subobject criterion recovers exactly $M$. Filtered colimits of nonempty sets are nonempty, so this is an L-accessible class. Its full subcategory has no initial object: a singleton has more than one map to a two-element set, and a larger set does not have a unique map either. It is therefore accessible but not cocomplete.

Under the hypotheses of Proposition~5.7, the statement independent of a chosen filtration is the existence of an essentially small family whose filtered-colimit closure is the given class. To pass from a small family to a bounded one, bound the filtration indices of its representatives. To pass back, take representatives of the bounded level. At a sufficiently large cardinal satisfying the accessibility requirement, Proposition~5.7 identifies this closure with the class detected by cofinally many good small subobjects. The cardinal filtration is then a means of making the construction uniform; the resulting notion of L-accessibility does not retain that choice.

There is a final scope convention to keep visible. Definition~5.12 restates $(a),(b),(c),(d)$ but omits $(e)$. The definition of L-$\pi$-accessibility itself makes sense with those hypotheses. Its cited classification through Proposition~5.7, however, was proved using $(e)$, and our small-subobject argument retains that hypothesis. We have not supplied a proof removing it from the classification, just as we have not removed the rank restriction from Proposition~5.9. These qualifications concern the claimed equivalences under weaker hypotheses, not the definition or the preceding example.

This is the part of the chapter which is closest to the later theory of accessible and accessibly embedded subcategories. It supplies a method for obtaining a small amount of generating data, while the earlier factorization argument turns small generating \emph{maps} into lifting factorizations. These are different passages: existence of a small generating family under filtered colimits does not by itself assert that the same maps generate a weak factorization system. The additional closure and lifting arguments remain necessary.
\end{dcomment}

